% Ch12 self-study -- "I do / You do" ladder for Zumdahl chapter 12.
% Twelve ladders, 48 questions, sixteen figures.
%
% SCOPE: nothing here is skipped. TEXTBOOK-MAP maps 12.1-12.7 onto CED
% 5.1 through 5.11 -- the whole of Unit 5 -- and the skip list contains
% no part of this chapter. It is the first chapter in the self-study
% series that is assessed end to end.
%
% Unit conventions in this file: rate constants are written in math
% mode ($\mathrm{M^{-1}\,s^{-1}}$) rather than through siunitx, because
% \per\Molar\squared ordering is a trap and k units change with order.
\documentclass{apchem-worksheet}

\apcunitword{Chapter}
\apcunit{12}
\apcunittitle{Chemical Kinetics}
\apcdoctitle{Self-Study \textbullet\ Chapter 12, I~do / You~do}
\apcsubtitle{Zumdahl \S12.1--12.7 \textbullet\ PDF pp.\ 594--632}

\begin{document}
\apctitleblock

\begin{apcnote}[How to use these notes --- read this first]
Each skill is a \textbf{ladder}: a fully worked example in the
solid-framed box, then \textbf{four} YOUR TURN questions in the dashed
box --- same skill, new numbers, no help. Work all four before checking
the gray \emph{check:} line; a tracker tick needs all four right first
time.

\textbf{Nothing in this chapter is skipped.} Every section, \S12.1
through \S12.7, is assessed --- they map onto CED topics 5.1 to 5.11,
which is the whole of Unit 5. This is the first chapter in the series
you have to learn end to end.

\textbf{The one idea the whole chapter turns on.} Thermodynamics tells
you \emph{whether} a reaction can happen. Kinetics tells you
\emph{how fast}, and the two are independent. Diamond turning into
graphite is thermodynamically downhill and yet takes longer than the age
of the Earth --- because the rate, not the favourability, is what you
observe.

\textbf{The habit that prevents most lost marks.} Orders come from
\textbf{experiment}, never from the balanced equation. Write that on the
inside of your eyelids. The single commonest error in Unit 5 is reading
a rate law off the coefficients of an overall equation.
\end{apcnote}

% =====================================================================
\section*{\sffamily Ladder 1 \textbullet\ Reaction rate and
stoichiometry}
\zsec{12.1}

Rate is a change in concentration per unit time. Because the species in
a reaction are consumed and produced at different speeds, we divide by
the coefficient so that \emph{one} number describes the reaction.

For $a\ce{A} + b\ce{B} \longrightarrow c\ce{C} + d\ce{D}$:
\[ \text{rate}
   = -\frac{1}{a}\frac{\Delta[\ce{A}]}{\Delta t}
   = -\frac{1}{b}\frac{\Delta[\ce{B}]}{\Delta t}
   = +\frac{1}{c}\frac{\Delta[\ce{C}]}{\Delta t}
   = +\frac{1}{d}\frac{\Delta[\ce{D}]}{\Delta t} \]

\begin{center}
\begin{tikzpicture}
\begin{axis}[
  width=10.5cm, height=6.2cm,
  xlabel={time (s)}, ylabel={concentration (\si{\Molar})},
  xmin=0, xmax=100, ymin=0, ymax=1.05,
  xtick={0,20,40,60,80,100}, ytick={0,0.2,0.4,0.6,0.8,1.0},
  grid=major, grid style={apcgray!25},
  label style={font=\footnotesize}, tick label style={font=\footnotesize},
]
% No legend: 'east' is not a valid pgfplots legend pos, and every corner
% of this plot already has a curve running through it. The two curves
% are labelled inline instead, in the empty band between them.
\addplot[seriesA, thick, domain=0:100, samples=80] {exp(-x/32)};
\addplot[seriesB, thick, domain=0:100, samples=80] {1-exp(-x/32)};
% tangent at t = 30 to the decay curve: y = e^{-30/32} = 0.3916,
% slope = -0.3916/32 = -0.01224
\addplot[black, dashed, domain=6:54, samples=2] {0.3916-0.01224*(x-30)};
\addplot[only marks, mark=*, mark size=1.6pt, black]
  coordinates {(30,0.3916)};
\node[font=\footnotesize,anchor=west] at (axis cs:56,0.70)
  {product --- \textbf{rises}};
\node[font=\footnotesize,anchor=west] at (axis cs:56,0.32)
  {reactant --- \textbf{falls}};
% Anchored WEST at x=24, in the empty band above both curves.  Anchored
% east at 29 these ran off the left edge and were clipped to
% "e of the tangent" / "stantaneous rate".
\node[font=\footnotesize,anchor=west] at (axis cs:24,0.97)
  {slope of the \textbf{tangent}};
\node[font=\footnotesize,anchor=west] at (axis cs:24,0.87)
  {$=$ instantaneous rate};
\end{axis}
\end{tikzpicture}
\end{center}

\begin{apctrap}
\textbf{Rate is always reported as a positive number.} A reactant's
concentration falls, so $\Delta[\ce{A}]/\Delta t$ is negative --- the
minus sign in the definition is there to cancel that, not to make the
rate negative. An answer of $\SI{-2.0e-3}{\Molar\per\second}$ for a rate
is wrong on sight.
\end{apctrap}

\begin{workedexample}[I do: one rate, four different species]
\textbf{For $\ce{2N2O5 -> 4NO2 + O2}$, \ce{N2O5} disappears at
\SI{4.0e-3}{\Molar\per\second}. Find the reaction rate and the rate of
appearance of each product.}

\textbf{Step 1 --- get the reaction rate.} Divide the disappearance of
\ce{N2O5} by its coefficient, 2:
\[ \text{rate} = \frac{1}{2}\left(\SI{4.0e-3}{\Molar\per\second}\right)
   = \SI{2.0e-3}{\Molar\per\second} \]

\textbf{Step 2 --- multiply back out for each species.} Each species
changes at (rate $\times$ its own coefficient):
\[ \frac{\Delta[\ce{NO2}]}{\Delta t}
   = 4 \times \SI{2.0e-3}{} = \SI{8.0e-3}{\Molar\per\second} \]
\[ \frac{\Delta[\ce{O2}]}{\Delta t}
   = 1 \times \SI{2.0e-3}{} = \SI{2.0e-3}{\Molar\per\second} \]

\textbf{Sanity-check against the chemistry.} Four \ce{NO2} appear for
every two \ce{N2O5} consumed, so \ce{NO2} should appear twice as fast as
\ce{N2O5} disappears --- \SI{8.0e-3}{} against \SI{4.0e-3}{}. It does.
That ratio check catches almost every slip in this ladder.

\textbf{Average versus instantaneous.} Concentration-versus-time curves
are curved, so the rate is not constant. A rate taken between two times
is an \textbf{average} rate (the slope of a straight line joining two
points); a rate at one instant is the slope of the \textbf{tangent}
there. Unless a question says otherwise, ``the rate'' means the
instantaneous rate, and the \textbf{initial rate} --- the tangent at
$t = 0$ --- is the one Ladder 3 depends on.
\end{workedexample}

\begin{yourturn}[YOUR TURN 1 --- four questions]
For $\ce{2NO + O2 -> 2NO2}$, \ce{O2} disappears at
\SI{0.020}{\Molar\per\second}.
\begin{enumerate}[leftmargin=*,itemsep=0.5em]
  \item What is the reaction rate? \workspace[1.4cm]
  \item How fast does \ce{NO} disappear? \workspace[1.4cm]
  \item How fast does \ce{NO2} appear? \workspace[1.4cm]
  \item Why is there a minus sign in front of $\Delta[\ce{A}]/\Delta t$
        for a reactant? \workspace[1.6cm]
\end{enumerate}
\selfcheck{(a) \SI{0.020}{\Molar\per\second} \quad (b)
\SI{0.040}{\Molar\per\second} \quad (c)
\SI{0.040}{\Molar\per\second} \quad (d) to make the rate come out
positive}
\keyonly{\begin{apcanswer}(a) \ce{O2} has coefficient 1, so the rate is
$\mathbf{\SI{0.020}{\Molar\per\second}}$. (b) \ce{NO} has coefficient 2,
so it disappears at $2 \times 0.020 =
\mathbf{\SI{0.040}{\Molar\per\second}}$. (c) \ce{NO2} also has
coefficient 2, so it appears at
$\mathbf{\SI{0.040}{\Molar\per\second}}$. (d) Because a reactant's
concentration \textbf{decreases}, so $\Delta[\ce{A}]$ is negative and
$\Delta[\ce{A}]/\Delta t$ is negative. Rate is defined as a positive
quantity, so the minus sign cancels that --- it is a bookkeeping device,
not a statement that anything is negative.\end{apcanswer}}
\end{yourturn}

% =====================================================================
\section*{\sffamily Ladder 2 \textbullet\ Rate laws and reaction order}
\zsec{12.2--12.3}

A rate law says how the rate depends on concentration:
\[ \text{rate} = k[\ce{A}]^{m}[\ce{B}]^{n} \]
$m$ and $n$ are the \textbf{orders}; $m + n$ is the \textbf{overall
order}; $k$ is the \textbf{rate constant}.

\begin{center}
\begin{tikzpicture}
\begin{axis}[
  width=10.6cm, height=5.4cm,
  xlabel={$[\ce{A}]$}, ylabel={rate},
  xmin=0, xmax=1.05, ymin=0, ymax=1.15,
  xtick=\empty, ytick=\empty,
  axis lines=left,
  legend pos=north west, legend style={font=\footnotesize,
    cells={anchor=west}},
  label style={font=\footnotesize},
]
\addplot[seriesA, thick, domain=0:1, samples=2] {0.55};
\addlegendentry{zero order --- rate independent of $[\ce{A}]$}
\addplot[seriesB, thick, domain=0:1, samples=2] {x};
\addlegendentry{first order --- straight line through the origin}
\addplot[seriesC, thick, domain=0:1, samples=60] {x^2};
\addlegendentry{second order --- upward curve}
\end{axis}
\end{tikzpicture}
\end{center}

\begin{apctrap}
\textbf{Orders are found by experiment. They are not the coefficients.}
For $\ce{2N2O5 -> 4NO2 + O2}$ the measured rate law is
rate $= k[\ce{N2O5}]$ --- \emph{first} order, not second, despite the
coefficient of 2. The only time coefficients give orders is for a single
\textbf{elementary step} (Ladder 7), and an overall balanced equation is
almost never one.
\end{apctrap}

\begin{workedexample}[I do: reading a rate law]
\textbf{A reaction has rate $= k[\ce{A}]^{2}[\ce{B}]$. Answer four
questions about it.}

\textbf{What is the order in each reactant?} Second order in \ce{A},
first order in \ce{B}. Read the exponents; a missing exponent means 1.

\textbf{What is the overall order?} $2 + 1 = \mathbf{3}$. Add the
exponents.

\textbf{What happens if $[\ce{A}]$ is doubled?} Rate changes by
$2^{2} = \mathbf{4}$ --- it quadruples. Raise the factor to the power of
the order.

\textbf{What happens if $[\ce{B}]$ is tripled?} Rate changes by
$3^{1} = \mathbf{3}$ --- it triples.

\textbf{And both at once?} Multiply the effects: $4 \times 3 = 12$.
Doubling \ce{A} and tripling \ce{B} makes the reaction twelve times
faster.

\textbf{A zero-order reactant is worth understanding, not memorising.}
If rate $= k[\ce{A}]^{0}[\ce{B}] = k[\ce{B}]$, then changing
$[\ce{A}]$ does nothing at all. That is not a trick --- it usually means
\ce{A} is not involved in the slow step, or that a surface or enzyme is
already saturated with it, so adding more has nowhere to go.
\end{workedexample}

\begin{yourturn}[YOUR TURN 2 --- four questions]
A reaction has rate $= k[\ce{X}][\ce{Y}]^{2}$.
\begin{enumerate}[leftmargin=*,itemsep=0.5em]
  \item What is the overall order? \workspace[1.4cm]
  \item By what factor does the rate change if $[\ce{Y}]$ is doubled?
        \workspace[1.4cm]
  \item By what factor does the rate change if $[\ce{X}]$ is halved?
        \workspace[1.4cm]
  \item A different reaction is $\ce{2A + B -> C}$ with measured rate
        $= k[\ce{A}][\ce{B}]$. Is the order in \ce{A} equal to 2?
        Explain. \workspace[1.8cm]
\end{enumerate}
\selfcheck{(a) 3 \quad (b) $4\times$ \quad (c) halves \quad (d) no ---
orders come from experiment}
\keyonly{\begin{apcanswer}(a) $1 + 2 = \mathbf{3}$. (b)
$2^{2} = \mathbf{4}$ --- the rate quadruples. (c)
$(1/2)^{1} = \mathbf{1/2}$ --- the rate \textbf{halves}. (d)
\textbf{No.} The measured order in \ce{A} is \textbf{1}, because the
rate law says so. The coefficient 2 in the balanced equation is
irrelevant to the rate law --- orders are determined experimentally and
match coefficients only for an elementary step, which an overall
equation almost never is.\end{apcanswer}}
\end{yourturn}

% =====================================================================
\section*{\sffamily Ladder 3 \textbullet\ The method of initial rates}
\zsec{12.2--12.3}

The standard experimental route to a rate law. Change \emph{one}
concentration at a time and watch what the rate does.

\begin{center}
\begin{tikzpicture}[font=\sffamily\footnotesize]
  \node[font=\sffamily\small\bfseries] at (5.2,2.6)
    {compare two experiments that differ in only one concentration};
  \node[anchor=west,align=left] at (0.2,1.9)
    {concentration changes by a factor of $f$};
  \node[anchor=west,align=left] at (0.2,1.35)
    {rate changes by a factor of $r$};
  \draw[apcgray] (0.1,0.95) -- (10.3,0.95);
  % The examples get their OWN line.  Both of these used to sit at
  % y=0.35, one anchored west and one east, and they printed straight
  % through each other in the middle of the figure.
  \node[anchor=west,font=\sffamily\bfseries] at (0.2,0.5)
    {then the order is the power that solves $f^{\,\text{order}} = r$};
  \node[anchor=west,font=\sffamily\footnotesize\itshape] at (0.2,-0.05)
    {$2 \to 2$: order 1 \textbullet\ $2 \to 4$: order 2 \textbullet\
     $2 \to 1$: order 0};
\end{tikzpicture}
\end{center}

\begin{workedexample}[I do: extract a rate law from a data table]
\textbf{For $\ce{A + B -> products}$:}

\begin{center}
\renewcommand{\arraystretch}{1.2}
\begin{tabular}{@{}cccc@{}}
\toprule
\textbf{Exp} & $[\ce{A}]_{0}$ (\si{\Molar}) & $[\ce{B}]_{0}$
  (\si{\Molar}) & \textbf{initial rate}
  (\si{\Molar\per\second}) \\
\midrule
1 & 0.10 & 0.10 & \num{2.0e-3} \\
2 & 0.20 & 0.10 & \num{8.0e-3} \\
3 & 0.10 & 0.20 & \num{4.0e-3} \\
\bottomrule
\end{tabular}
\end{center}

\textbf{Order in \ce{A} --- use experiments 1 and 2}, because
$[\ce{B}]$ is the same in both. $[\ce{A}]$ doubles
($0.10 \to 0.20$) and the rate goes from \num{2.0e-3} to \num{8.0e-3},
a factor of 4.
\[ 2^{m} = 4 \;\Longrightarrow\; m = \mathbf{2} \]

\textbf{Order in \ce{B} --- use experiments 1 and 3}, where $[\ce{A}]$
is held fixed. $[\ce{B}]$ doubles and the rate doubles
(\num{2.0e-3} to \num{4.0e-3}).
\[ 2^{n} = 2 \;\Longrightarrow\; n = \mathbf{1} \]

\textbf{So the rate law is} $\text{rate} = k[\ce{A}]^{2}[\ce{B}]$,
overall order 3.

\textbf{Now find $k$} by substituting any one experiment --- use
experiment 1:
\[ k = \frac{\text{rate}}{[\ce{A}]^{2}[\ce{B}]}
     = \frac{\num{2.0e-3}}{(0.10)^{2}(0.10)}
     = \frac{\num{2.0e-3}}{\num{1.0e-3}}
     = 2.0\;\mathrm{M^{-2}\,s^{-1}} \]

\textbf{Then check $k$ against a different row} --- this is free and it
catches errors. Experiment 2:
$2.0 \times (0.20)^{2} \times 0.10 = 2.0 \times 0.040 \times 0.10 =
\num{8.0e-3}$. Matches. If $k$ does not come out the same from every
row, an order is wrong.
\end{workedexample}

\begin{yourturn}[YOUR TURN 3 --- four questions]
For $\ce{X + Y -> products}$:
\begin{center}
\renewcommand{\arraystretch}{1.15}
\begin{tabular}{@{}cccc@{}}
\toprule
\textbf{Exp} & $[\ce{X}]_{0}$ (\si{\Molar}) & $[\ce{Y}]_{0}$
  (\si{\Molar}) & \textbf{initial rate} (\si{\Molar\per\second}) \\
\midrule
1 & 0.10 & 0.10 & \num{1.5e-3} \\
2 & 0.20 & 0.10 & \num{3.0e-3} \\
3 & 0.10 & 0.20 & \num{1.5e-3} \\
\bottomrule
\end{tabular}
\end{center}
\begin{enumerate}[leftmargin=*,itemsep=0.5em]
  \item Find the order in \ce{X}. \workspace[1.6cm]
  \item Find the order in \ce{Y}. \workspace[1.6cm]
  \item Write the rate law and give the overall order.
        \workspace[1.6cm]
  \item Calculate $k$, with units. \workspace[1.8cm]
\end{enumerate}
\selfcheck{(a) 1 \quad (b) 0 \quad (c) rate $= k[\ce{X}]$, first order
\quad (d) \SI{1.5e-2}{\per\second}}
\keyonly{\begin{apcanswer}(a) Experiments 1 and 2 hold $[\ce{Y}]$
fixed. $[\ce{X}]$ doubles and the rate doubles, so
$2^{m} = 2$ and $m = \mathbf{1}$. (b) Experiments 1 and 3 hold
$[\ce{X}]$ fixed. $[\ce{Y}]$ doubles and the rate does
\textbf{not change}, so $2^{n} = 1$ and $n = \mathbf{0}$. (c)
$\text{rate} = k[\ce{X}]$, overall order \textbf{1} --- \ce{Y} does not
appear, because raising it changes nothing. (d) $k =
\text{rate}/[\ce{X}] = \num{1.5e-3}/0.10 =
\mathbf{\SI{1.5e-2}{\per\second}}$. A first-order rate constant always
has units of $\mathrm{s^{-1}}$.\end{apcanswer}}
\end{yourturn}

% =====================================================================
\section*{\sffamily Ladder 4 \textbullet\ The units of $k$}
\zsec{12.2}

The units of $k$ are not fixed --- they are whatever makes the rate law
come out in \si{\Molar\per\second}. That makes them a free check on
your order.

\begin{center}
\begin{tikzpicture}[font=\sffamily\footnotesize]
  \node[font=\sffamily\small\bfseries] at (5.2,2.9)
    {units of $k$ for overall order $n$};
  \foreach \x/\o/\u in {%
      0.6/{0}/{$\mathrm{M\,s^{-1}}$},
      3.2/{1}/{$\mathrm{s^{-1}}$},
      5.8/{2}/{$\mathrm{M^{-1}\,s^{-1}}$},
      8.4/{3}/{$\mathrm{M^{-2}\,s^{-1}}$}} {
    \node[font=\sffamily\bfseries] at (\x,2.25) {order \o};
    \node at (\x,1.65) {\u};}
  \draw[apcgray] (0.1,1.2) -- (10.3,1.2);
  \node[anchor=west] at (0.2,0.75)
    {the pattern: $\mathrm{M^{\,1-n}\,s^{-1}}$ --- each extra order
     removes one power of M};
  \node[anchor=west,font=\sffamily\itshape] at (0.2,0.25)
    {so if you are told the units of $k$, you already know the overall
     order};
\end{tikzpicture}
\end{center}

\begin{workedexample}[I do: derive the units instead of memorising them]
\textbf{Never memorise this table. Derive it in five seconds.}

Rate always has units \si{\Molar\per\second}, i.e.\
$\mathrm{M\,s^{-1}}$. For a second-order reaction,
rate $= k[\ce{A}]^{2}$, so
\[ k = \frac{\text{rate}}{[\ce{A}]^{2}}
     = \frac{\mathrm{M\,s^{-1}}}{\mathrm{M^{2}}}
     = \mathrm{M^{-1}\,s^{-1}} \]

\textbf{The same move for any order.} Divide $\mathrm{M\,s^{-1}}$ by
$\mathrm{M}^{n}$, giving $\mathrm{M^{\,1-n}\,s^{-1}}$. For $n=0$ that is
$\mathrm{M\,s^{-1}}$; for $n=1$, $\mathrm{s^{-1}}$; for $n=3$,
$\mathrm{M^{-2}\,s^{-1}}$.

\textbf{Now use it backwards, which is where the marks are.} If a
question hands you $k = 0.045\;\mathrm{M^{-1}\,s^{-1}}$, you already
know the reaction is \textbf{second order overall} before reading
anything else. Examiners use this to test whether you understand the
relationship or merely memorised a table.

\textbf{And use it as a check.} If you determine orders from a data
table and then compute a $k$ whose units disagree with your overall
order, one of the two is wrong --- go back before continuing.
\end{workedexample}

\begin{yourturn}[YOUR TURN 4 --- four questions]
\begin{enumerate}[leftmargin=*,itemsep=0.5em]
  \item Give the units of $k$ for a first-order reaction.
        \workspace[1.4cm]
  \item $k = 0.020\;\mathrm{M^{-1}\,s^{-1}}$. What is the overall
        order? \workspace[1.4cm]
  \item $k = 3.5\;\mathrm{M\,s^{-1}}$. What is the overall order?
        \workspace[1.4cm]
  \item Derive the units of $k$ for a reaction that is first order in
        \ce{A} and second order in \ce{B}. \workspace[1.8cm]
\end{enumerate}
\selfcheck{(a) $\mathrm{s^{-1}}$ \quad (b) 2 \quad (c) 0 \quad (d)
$\mathrm{M^{-2}\,s^{-1}}$}
\keyonly{\begin{apcanswer}(a) $\mathbf{\mathrm{s^{-1}}}$. (b)
\textbf{Second order} --- $\mathrm{M^{-1}}$ means $1-n = -1$, so
$n = 2$. (c) \textbf{Zero order} --- $\mathrm{M^{1}}$ means $1-n = 1$,
so $n = 0$, and the rate does not depend on concentration at all.
(d) Overall order is $1 + 2 = 3$, so $k = \mathrm{M\,s^{-1}} /
\mathrm{M^{3}} = \mathbf{\mathrm{M^{-2}\,s^{-1}}}$.\end{apcanswer}}
\end{yourturn}

% =====================================================================
\section*{\sffamily Ladder 5 \textbullet\ Integrated rate laws}
\zsec{12.4}

A rate law relates rate to concentration. An \emph{integrated} rate law
relates concentration to \textbf{time} --- which is what you actually
need to answer ``how much is left after ten minutes?''

\begin{center}
\renewcommand{\arraystretch}{1.5}
\begin{tabular}{@{}llll@{}}
\toprule
\textbf{Order} & \textbf{Integrated law} &
  \textbf{Linear plot} & \textbf{Slope} \\
\midrule
0 & $[\ce{A}] = -kt + [\ce{A}]_{0}$ &
  $[\ce{A}]$ vs $t$ & $-k$ \\
1 & $\ln[\ce{A}] = -kt + \ln[\ce{A}]_{0}$ &
  $\ln[\ce{A}]$ vs $t$ & $-k$ \\
2 & $\dfrac{1}{[\ce{A}]} = kt + \dfrac{1}{[\ce{A}]_{0}}$ &
  $1/[\ce{A}]$ vs $t$ & $+k$ \\
\bottomrule
\end{tabular}
\end{center}

\begin{center}
\begin{tikzpicture}[font=\sffamily\footnotesize]
  \foreach \i/\ttl/\yl in {%
      0/{zero order}/{$[\ce{A}]$},
      3.9/{first order}/{$\ln[\ce{A}]$},
      7.8/{second order}/{$1/[\ce{A}]$}} {
    \begin{scope}[xshift=\i cm]
      \node[font=\sffamily\small\bfseries] at (1.5,3.15) {\ttl};
      \draw[-{Stealth[length=2.2mm]},thick] (0.35,0.4) -- (0.35,2.65);
      \draw[-{Stealth[length=2.2mm]},thick] (0.35,0.4) -- (2.85,0.4);
      \node[anchor=east,font=\footnotesize] at (0.25,1.55) {\yl};
      \node[anchor=north,font=\footnotesize] at (1.6,0.3) {$t$};
    \end{scope}}
  % straight lines: 0 and 1 fall, 2 rises
  \draw[very thick] (0.55,2.35) -- (2.65,0.65);
  \begin{scope}[xshift=3.9cm]
    \draw[very thick] (0.55,2.35) -- (2.65,0.65);
  \end{scope}
  \begin{scope}[xshift=7.8cm]
    \draw[very thick] (0.55,0.65) -- (2.65,2.35);
  \end{scope}
  \node[font=\sffamily\footnotesize] at (5.3,-0.25)
    {plot all three; \textbf{whichever comes out straight} tells you the
     order};
\end{tikzpicture}
\end{center}

\begin{apcnote}[Why this is really a graphing skill]
On the exam you are far more likely to be \emph{given} data and asked
which order it is than to be told the order outright. The method is
always the same: make the three plots (or inspect three columns of
numbers) and find the one that is linear. A straight
$\ln[\ce{A}]$-versus-$t$ plot \textbf{is} the evidence for first order,
and saying so is what earns the mark --- not asserting the order and
moving on.
\end{apcnote}

\begin{workedexample}[I do: a first-order calculation]
\textbf{A first-order reaction has $k = \SI{0.0250}{\per\second}$ and
$[\ce{A}]_{0} = \SI{0.800}{\Molar}$. Find $[\ce{A}]$ after
\SI{60.0}{\second}.}

\textbf{Use the first-order integrated law} in its exponential form,
which is easier than the logarithmic one when you want a concentration:
\[ [\ce{A}] = [\ce{A}]_{0}\,e^{-kt} \]
\[ [\ce{A}] = 0.800 \times e^{-(0.0250)(60.0)}
            = 0.800 \times e^{-1.500}
            = 0.800 \times 0.2231
            = \mathbf{\SI{0.179}{\Molar}} \]

\textbf{Check it with half-lives, which is quicker than it looks.}
\[ t_{1/2} = \frac{0.693}{k} = \frac{0.693}{0.0250}
           = \SI{27.7}{\second} \]
\SI{60.0}{\second} is about \textbf{2.2} half-lives --- so a little more
than two, and roughly a quarter of the material should be left, with a
bit more decay on top. $0.800/4 = 0.200$, and our answer of
\SI{0.179}{\Molar} sits just below it. Consistent.

\textbf{That cross-check is worth doing every time.} It costs one line
and catches sign errors in the exponent, which are the commonest failure
here. A positive exponent would have given \SI{3.58}{\Molar} --- more
than you started with, and obviously impossible.
\end{workedexample}

\begin{yourturn}[YOUR TURN 5 --- four questions]
\begin{enumerate}[leftmargin=*,itemsep=0.5em]
  \item Which plot is linear for a second-order reaction?
        \workspace[1.4cm]
  \item A plot of $\ln[\ce{A}]$ against $t$ is a straight line with
        slope $\SI{-0.0400}{\per\second}$. Give the order and $k$.
        \workspace[1.8cm]
  \item For a first-order reaction with $k = \SI{0.100}{\per\second}$
        and $[\ce{A}]_{0} = \SI{1.00}{\Molar}$, find $[\ce{A}]$ after
        \SI{10.0}{\second}. \workspace[1.8cm]
  \item You are given concentration-versus-time data and no order.
        Describe how you would find the order. \workspace[1.8cm]
\end{enumerate}
\selfcheck{(a) $1/[\ce{A}]$ vs $t$ \quad (b) first order,
$k = \SI{0.0400}{\per\second}$ \quad (c) \SI{0.368}{\Molar} \quad (d)
plot all three, find the straight one}
\keyonly{\begin{apcanswer}(a) $\mathbf{1/[\ce{A}]}$ \textbf{against}
$t$, with slope $+k$. (b) A linear $\ln[\ce{A}]$ plot means
\textbf{first order}; the slope is $-k$, so
$k = \mathbf{\SI{0.0400}{\per\second}}$ (positive). (c)
$[\ce{A}] = 1.00 \times e^{-(0.100)(10.0)} = e^{-1.00} =
\mathbf{\SI{0.368}{\Molar}}$. (d) Make \textbf{all three} plots ---
$[\ce{A}]$, $\ln[\ce{A}]$ and $1/[\ce{A}]$, each against $t$ --- and see
which one is a straight line. That plot identifies the order, and its
slope gives $k$ ($-k$ for zero and first order, $+k$ for
second).\end{apcanswer}}
\end{yourturn}

% =====================================================================
\section*{\sffamily Ladder 6 \textbullet\ Half-life}
\zsec{12.4}

The half-life $t_{1/2}$ is the time for half the reactant to be
consumed. Its behaviour differs sharply between orders, and that
difference is itself a way to identify the order.

\begin{center}
\begin{tikzpicture}
\begin{axis}[
  width=10.5cm, height=5.8cm,
  xlabel={time (s)}, ylabel={$[\ce{A}]$ (\si{\Molar})},
  xmin=0, xmax=120, ymin=0, ymax=1.1,
  xtick={0,30,60,90,120}, ytick={0,0.25,0.5,0.75,1.0},
  grid=major, grid style={apcgray!25},
  label style={font=\footnotesize}, tick label style={font=\footnotesize},
]
\addplot[seriesA, very thick, domain=0:120, samples=90] {exp(-0.0231*x)};
\foreach \t/\c in {30/0.5, 60/0.25, 90/0.125} {
  \addplot[black, dotted, thick, domain=0:\t, samples=2] {\c};
}
\draw[dotted, thick] (axis cs:30,0) -- (axis cs:30,0.5);
\draw[dotted, thick] (axis cs:60,0) -- (axis cs:60,0.25);
\draw[dotted, thick] (axis cs:90,0) -- (axis cs:90,0.125);
\node[font=\footnotesize,anchor=west] at (axis cs:44,0.86)
  {each drop by half takes the \textbf{same} \SI{30}{\second}};
\node[font=\footnotesize,anchor=west] at (axis cs:44,0.74)
  {--- the signature of \textbf{first order}};
\end{axis}
\end{tikzpicture}
\end{center}

\begin{center}
\begin{tikzpicture}[font=\sffamily\footnotesize]
  \node[font=\sffamily\small\bfseries] at (5.3,2.55)
    {how successive half-lives behave};
  \foreach \y/\o/\f/\b in {%
      1.85/{zero order}/{$t_{1/2} = \dfrac{[\ce{A}]_{0}}{2k}$}/{each one
        is \textbf{shorter}},
      1.0/{first order}/{$t_{1/2} = \dfrac{0.693}{k}$}/{all
        \textbf{the same}},
      0.15/{second order}/{$t_{1/2} = \dfrac{1}{k[\ce{A}]_{0}}$}/{each
        one is \textbf{longer}}} {
    % formula column at 3.0, not 2.5: "second order" in bold sans is
    % wider than 2.3cm and touched the formula beside it.
    \node[anchor=west,font=\sffamily\bfseries] at (0.2,\y) {\o};
    \node[anchor=west] at (3.0,\y) {\f};
    \node[anchor=west] at (7.0,\y) {\b};}
  \draw[apcgray] (0.1,2.3) -- (10.4,2.3);
\end{tikzpicture}
\end{center}

\begin{apctrap}
\textbf{Only the first-order half-life is independent of
concentration.} That is why $t_{1/2} = 0.693/k$ contains no
$[\ce{A}]_{0}$ --- and why radioactive decay, which is always first
order, has a fixed half-life you can quote for an isotope. Quoting a
single half-life for a second-order reaction is meaningless, because it
doubles every time round.
\end{apctrap}

\begin{workedexample}[I do: half-lives without a calculator]
\textbf{A first-order reaction has $t_{1/2} = \SI{25}{\second}$. What
fraction of the reactant remains after \SI{100}{\second}?}

\textbf{Count half-lives rather than using the exponential.}
\[ \frac{\SI{100}{\second}}{\SI{25}{\second}} = 4
   \text{ half-lives} \]
Each one leaves half of what was there, so after four:
\[ \left(\tfrac{1}{2}\right)^{4} = \tfrac{1}{16}
   = \mathbf{0.0625} = 6.25\% \]

\textbf{This is exactly why MCQ sections can be no-calculator.} Whole
numbers of half-lives are meant to be done in your head: 1 half-life
leaves $1/2$, then $1/4$, $1/8$, $1/16$. If a question gives you a time
that is a neat multiple of $t_{1/2}$, it is telling you to count rather
than compute.

\textbf{Going the other way.} If a first-order reaction is 87.5\%
complete, then $12.5\% = 1/8$ remains, which is three half-lives.
Recognising $1/2$, $1/4$, $1/8$, $1/16$ as one, two, three and four
half-lives turns a calculation into a glance.

\textbf{And the link back to $k$:} $t_{1/2} = 0.693/k$, so a large $k$
means a short half-life. The two are just different ways of stating the
same speed.
\end{workedexample}

\begin{yourturn}[YOUR TURN 6 --- four questions]
\begin{enumerate}[leftmargin=*,itemsep=0.5em]
  \item A first-order reaction has $k = \SI{0.0350}{\per\second}$. Find
        $t_{1/2}$. \workspace[1.6cm]
  \item What fraction remains after 3 half-lives?
        \workspace[1.4cm]
  \item A first-order reaction is 75\% complete. How many half-lives
        have passed? \workspace[1.6cm]
  \item Successive half-lives of a reaction get longer each time. What
        order is it? \workspace[1.6cm]
\end{enumerate}
\selfcheck{(a) \SI{19.8}{\second} \quad (b) $1/8$ \quad (c) 2 \quad
(d) second order}
\keyonly{\begin{apcanswer}(a) $t_{1/2} = 0.693/0.0350 =
\mathbf{\SI{19.8}{\second}}$. (b) $(1/2)^{3} = \mathbf{1/8}$, or
12.5\%. (c) 75\% complete means \textbf{25\% remains}, and
$1/4 = (1/2)^{2}$, so \textbf{2} half-lives. (d) \textbf{Second order.}
Its half-life is $1/(k[\ce{A}]_{0})$, so as the concentration falls the
half-life rises --- it doubles each time. A first-order half-life would
stay constant and a zero-order one would
shorten.\end{apcanswer}}
\end{yourturn}

% =====================================================================
\section*{\sffamily Ladder 7 \textbullet\ Elementary steps and
molecularity}
\zsec{12.5}

An \textbf{elementary step} is a single collision event --- what
actually happens, not a summary. For an elementary step, and only for an
elementary step, the coefficients \emph{are} the orders.

\begin{center}
\begin{tikzpicture}[font=\sffamily\footnotesize]
  \foreach \i/\nm/\eq/\rl in {%
      0/{unimolecular}/{$\ce{A -> products}$}/{rate $= k[\ce{A}]$},
      3.8/{bimolecular}/{$\ce{A + B -> products}$}/{rate $=
        k[\ce{A}][\ce{B}]$},
      7.6/{termolecular}/{$\ce{A + B + C -> products}$}/{rate $=
        k[\ce{A}][\ce{B}][\ce{C}]$}} {
    \begin{scope}[xshift=\i cm]
      \node[font=\sffamily\small\bfseries] at (1.6,3.15) {\nm};
      \node[font=\footnotesize] at (1.6,1.05) {\eq};
      \node[font=\footnotesize] at (1.6,0.5) {\rl};
    \end{scope}}
  % art: 1 particle / 2 colliding / 3 colliding
  \fill[apcgray!65] (1.6,2.2) circle (0.2);
  \draw[-{Stealth[length=2mm]}] (1.6,1.85) -- (1.6,1.5);
  \begin{scope}[xshift=3.8cm]
    \fill[apcgray!65] (1.25,2.2) circle (0.2);
    \fill[apcgray!30] (1.95,2.2) circle (0.2);
    \draw[-{Stealth[length=2mm]}] (1.42,2.2) -- (1.62,2.2);
    \draw[{Stealth[length=2mm]}-] (1.58,2.2) -- (1.78,2.2);
  \end{scope}
  \begin{scope}[xshift=7.6cm]
    \fill[apcgray!65] (1.05,2.2) circle (0.19);
    \fill[apcgray!30] (1.6,2.2) circle (0.19);
    \fill[apcgray!45] (2.15,2.2) circle (0.19);
  \end{scope}
  % Note moved BELOW the rate law.  At y=1.75 its top line ran into the
  % three circles sitting at y=2.2.
  \node[font=\sffamily\footnotesize,align=center] at (9.2,-0.2)
    {\emph{rare} --- three particles meeting\\
     at once is highly improbable};
\end{tikzpicture}
\end{center}

\begin{workedexample}[I do: the one place coefficients give orders]
\textbf{Why can you read a rate law off an elementary step but not off
an overall equation?}

\textbf{Because an elementary step describes an actual collision.} If
the step is $\ce{NO2 + NO2 -> NO3 + NO}$, then the event is literally
two \ce{NO2} molecules meeting. Double $[\ce{NO2}]$ and each molecule
has twice as many partners available while there are twice as many
molecules looking --- so the collision frequency, and the rate, goes up
by $2 \times 2 = 4$. That is second order, and it follows from the
coefficient 2 because the coefficient is counting colliding particles.

\textbf{An overall equation summarises, and summaries lose
information.} $\ce{2N2O5 -> 4NO2 + O2}$ does not mean two \ce{N2O5}
molecules collide. It means that when the whole multi-step process is
tallied up, two are consumed. The measured rate law is
$k[\ce{N2O5}]$ --- first order --- because the slow step involves only
one molecule.

\textbf{Why termolecular steps are essentially never proposed.} Getting
two particles to meet with enough energy and the right orientation is
already demanding. Requiring a third to arrive at the same instant makes
the probability vanishingly small. If you write a mechanism with a
termolecular step, expect to be asked to justify it --- and you usually
cannot.

\textbf{The test to apply:} has the question told you the step is
elementary (or is it labelled ``slow''/``fast'' inside a mechanism)? If
yes, use the coefficients. If it is an overall equation, you need
experimental data.
\end{workedexample}

\begin{yourturn}[YOUR TURN 7 --- four questions]
\begin{enumerate}[leftmargin=*,itemsep=0.5em]
  \item Give the rate law for the elementary step
        $\ce{2NO2 -> N2O4}$. \workspace[1.6cm]
  \item Give the molecularity of the elementary step
        $\ce{O3 -> O2 + O}$. \workspace[1.4cm]
  \item Give the rate law for the elementary step
        $\ce{NO + O3 -> NO2 + O2}$. \workspace[1.6cm]
  \item Why are termolecular steps rare? \workspace[1.8cm]
\end{enumerate}
\selfcheck{(a) rate $= k[\ce{NO2}]^{2}$ \quad (b) unimolecular \quad
(c) rate $= k[\ce{NO}][\ce{O3}]$ \quad (d) three-body collisions are
improbable}
\keyonly{\begin{apcanswer}(a) $\text{rate} =
\mathbf{k[\ce{NO2}]^{2}}$ --- for an elementary step the coefficient
becomes the order. (b) \textbf{Unimolecular}: one particle decomposes
on its own. (c) $\text{rate} = \mathbf{k[\ce{NO}][\ce{O3}]}$,
bimolecular and first order in each. (d) Because they require
\textbf{three particles to collide simultaneously} with sufficient
energy and correct orientation. Two-body collisions are already
relatively infrequent; the chance of a third arriving at the same
instant is so small that termolecular steps are almost never a
plausible part of a mechanism.\end{apcanswer}}
\end{yourturn}

% =====================================================================
\section*{\sffamily Ladder 8 \textbullet\ Mechanisms and the
rate-determining step}
\zsec{12.5}

A mechanism is the list of elementary steps. The \textbf{slowest} step
sets the pace for the whole reaction, exactly as the narrowest point on
a road sets the traffic speed.

\begin{center}
\begin{tikzpicture}[font=\sffamily\footnotesize]
  \node[font=\sffamily\small\bfseries] at (5.3,3.5)
    {a valid mechanism must pass three tests};
  \foreach \y/\n/\t in {%
      2.75/{1}/{the steps must add up to the overall equation},
      2.1/{2}/{the rate law from the slow step must match the measured
        one},
      1.45/{3}/{every intermediate must be produced, then consumed}} {
    \node[anchor=west,font=\sffamily\bfseries] at (0.3,\y) {\n.};
    \node[anchor=west] at (0.85,\y) {\t};}
  \draw[apcgray] (0.2,1.05) -- (10.4,1.05);
  \node[anchor=west,align=left] at (0.3,0.5)
    {\textbf{intermediate}: made in an early step, used up later ---
     never in the overall equation};
  \node[anchor=west,align=left] at (0.3,0.0)
    {\textbf{catalyst}: used in an early step, \emph{regenerated}
     later --- also never in the overall equation};
\end{tikzpicture}
\end{center}

\begin{center}
\begin{tikzpicture}[font=\sffamily\footnotesize]
  \node[font=\sffamily\small\bfseries] at (5.0,2.95)
    {$\ce{NO2 + CO -> NO + CO2}$};
  % Equations start at 3.3, not 2.6 -- the bold sans label is wider than
  % 2.4cm and printed straight into the equation ("Step 1 (slow):NO2...").
  \node[anchor=west,font=\sffamily\bfseries] at (0.2,2.2)
    {Step 1 (\textbf{slow}):};
  \node[anchor=west] at (3.3,2.2) {$\ce{NO2 + NO2 -> NO3 + NO}$};
  \node[anchor=west,font=\sffamily\bfseries] at (0.2,1.55)
    {Step 2 (fast):};
  \node[anchor=west] at (3.3,1.55) {$\ce{NO3 + CO -> NO2 + CO2}$};
  \draw[apcgray] (0.2,1.15) -- (10.0,1.15);
  \node[anchor=west] at (0.2,0.7)
    {\ce{NO3} is made in step 1 and consumed in step 2 --- an
     \textbf{intermediate}};
  \node[anchor=west] at (0.2,0.15)
    {the slow step is bimolecular in \ce{NO2}, so
     $\text{rate} = k[\ce{NO2}]^{2}$};
\end{tikzpicture}
\end{center}

\begin{workedexample}[I do: check a mechanism, then get its rate law]
\textbf{Test the two-step mechanism above against all three
requirements.}

\textbf{Test 1 --- do the steps add to the overall equation?} Add them:
\[ \ce{2NO2 + NO3 + CO -> NO3 + NO + NO2 + CO2} \]
Cancel \ce{NO3} from both sides and one \ce{NO2} from both sides:
\[ \ce{NO2 + CO -> NO + CO2} \quad\checkmark \]

\textbf{Test 2 --- what rate law does it predict?} The slow step is
$\ce{NO2 + NO2 -> NO3 + NO}$, which is elementary, so its coefficients
give its orders:
\[ \text{rate} = k[\ce{NO2}]^{2} \]
This is the experimentally measured rate law for this reaction, so the
mechanism survives.

\textbf{Test 3 --- is the intermediate handled properly?} \ce{NO3} is
produced in step 1 and consumed in step 2, and it does not appear in the
overall equation. \quad\checkmark

\textbf{Now notice what the rate law does \emph{not} contain: \ce{CO}.}
Carbon monoxide is a reactant in the overall equation, yet the rate does
not depend on it at all --- because \ce{CO} appears only in the
\emph{fast} step, after the bottleneck. Adding more \ce{CO} cannot speed
up a process that is already waiting on step 1.

\textbf{And notice the order in \ce{NO2} is 2} while its coefficient in
the overall equation is 1. Both observations are the same lesson from
Ladder 2, now with a mechanism to explain \emph{why}.
\end{workedexample}

\begin{yourturn}[YOUR TURN 8 --- four questions]
Consider this mechanism for $\ce{2NO2 + F2 -> 2NO2F}$:
\begin{center}
\begin{tabular}{@{}ll@{}}
Step 1 (slow): & $\ce{NO2 + F2 -> NO2F + F}$ \\
Step 2 (fast): & $\ce{NO2 + F -> NO2F}$ \\
\end{tabular}
\end{center}
\begin{enumerate}[leftmargin=*,itemsep=0.5em]
  \item Show the steps add to the overall equation.
        \workspace[1.8cm]
  \item Identify the intermediate. \workspace[1.4cm]
  \item Write the predicted rate law. \workspace[1.6cm]
  \item Would doubling $[\ce{NO2}]$ double the rate? Explain.
        \workspace[1.8cm]
\end{enumerate}
\selfcheck{(a) they sum correctly \quad (b) \ce{F} atoms \quad (c) rate
$= k[\ce{NO2}][\ce{F2}]$ \quad (d) yes --- first order in \ce{NO2}}
\keyonly{\begin{apcanswer}(a) Adding gives $\ce{2NO2 + F2 + F -> 2NO2F +
F}$; cancelling the \ce{F} that appears on both sides leaves
$\ce{2NO2 + F2 -> 2NO2F}$. \checkmark\ (b) The \textbf{fluorine atom},
\ce{F} --- produced in step 1 and consumed in step 2, and absent from the
overall equation. (c) From the slow elementary step:
$\text{rate} = \mathbf{k[\ce{NO2}][\ce{F2}]}$. (d) \textbf{Yes.} The
rate law is first order in \ce{NO2}, so doubling it doubles the rate ---
even though \ce{NO2} has a coefficient of 2 in the overall equation. Only
one \ce{NO2} takes part in the rate-determining step; the second is
consumed later, in the fast step, where it cannot affect the
rate.\end{apcanswer}}
\end{yourturn}

% =====================================================================
\section*{\sffamily Ladder 9 \textbullet\ A fast first step:
pre-equilibrium}
\zsec{12.5}

When the \emph{first} step is fast and reversible, the rate-determining
step contains an intermediate --- and a rate law may not contain an
intermediate, because you cannot measure its concentration. This ladder
is how you get rid of it.

\begin{center}
\begin{tikzpicture}[font=\sffamily\footnotesize]
  \node[font=\sffamily\small\bfseries] at (5.0,3.15)
    {$\ce{2NO + O2 -> 2NO2}$};
  % Equations start at 4.9.  At 4.3 the long bold label overprinted the
  % equation and swallowed the leading coefficient, so "2NO <=> N2O2"
  % printed as "NO <=> N2O2" -- a chemistry error caused by a layout bug.
  \node[anchor=west,font=\sffamily\bfseries] at (0.2,2.45)
    {Step 1 (fast, \emph{reversible}):};
  \node[anchor=west] at (4.9,2.45)
    {$\ce{2NO <=> N2O2}$};
  \node[anchor=west,font=\sffamily\bfseries] at (0.2,1.8)
    {Step 2 (\textbf{slow}):};
  \node[anchor=west] at (4.9,1.8)
    {$\ce{N2O2 + O2 -> 2NO2}$};
  \draw[apcgray] (0.2,1.4) -- (10.4,1.4);
  \node[anchor=west] at (0.2,0.9)
    {slow step gives $\text{rate} = k_{2}[\ce{N2O2}][\ce{O2}]$ ---
     but \ce{N2O2} is an intermediate};
  \node[anchor=west] at (0.2,0.15)
    {step 1 at equilibrium gives
     $K = \dfrac{[\ce{N2O2}]}{[\ce{NO}]^{2}}$, so
     $[\ce{N2O2}] = K[\ce{NO}]^{2}$};
  \node[anchor=west,font=\sffamily\bfseries] at (0.2,-0.6)
    {substitute: $\text{rate} = k_{2}K[\ce{NO}]^{2}[\ce{O2}]
     = k[\ce{NO}]^{2}[\ce{O2}]$};
\end{tikzpicture}
\end{center}

\begin{workedexample}[I do: eliminating an intermediate]
\textbf{The rule that forces this whole procedure: a rate law may
contain only species whose concentrations you can control or measure ---
reactants, products, catalysts. Never an intermediate.}

\textbf{Start where the bottleneck is.} The slow step is
$\ce{N2O2 + O2 -> 2NO2}$, elementary, so
\[ \text{rate} = k_{2}[\ce{N2O2}][\ce{O2}] \]
This is correct but useless: \ce{N2O2} exists only fleetingly and you
cannot put a number on its concentration.

\textbf{Use the fast step to express it in measurable terms.} Step 1 is
fast and reversible, so it reaches equilibrium long before step 2 gets
anywhere. For $\ce{2NO <=> N2O2}$,
\[ K = \frac{[\ce{N2O2}]}{[\ce{NO}]^{2}}
   \quad\Longrightarrow\quad
   [\ce{N2O2}] = K[\ce{NO}]^{2} \]

\textbf{Substitute and fold the constants together.}
\[ \text{rate} = k_{2}K[\ce{NO}]^{2}[\ce{O2}]
               = k[\ce{NO}]^{2}[\ce{O2}] \]
where $k = k_{2}K$ is what an experiment actually measures. You cannot
separate $k_{2}$ from $K$ by measuring the rate --- only their product.

\textbf{Check against reality.} The experimentally observed rate law for
$\ce{2NO + O2 -> 2NO2}$ is third order overall, second in \ce{NO} and
first in \ce{O2}. The mechanism reproduces it exactly.

\textbf{And note the payoff.} A simple one-step collision of
$\ce{2NO + O2}$ would be \emph{termolecular} --- three particles at
once, which Ladder 7 said is implausible. The pre-equilibrium mechanism
gets the same third-order rate law out of two ordinary bimolecular
steps. That is why chemists proposed it.
\end{workedexample}

\begin{yourturn}[YOUR TURN 9 --- four questions]
\begin{enumerate}[leftmargin=*,itemsep=0.5em]
  \item Why may a rate law not contain an intermediate?
        \workspace[1.8cm]
  \item In the mechanism above, which species is the intermediate?
        \workspace[1.4cm]
  \item A mechanism has fast $\ce{A <=> B}$ then slow
        $\ce{B + C -> D}$. Write the rate law in terms of \ce{A} and
        \ce{C}. \workspace[1.8cm]
  \item Why is a termolecular one-step route for
        $\ce{2NO + O2 -> 2NO2}$ considered implausible?
        \workspace[1.8cm]
\end{enumerate}
\selfcheck{(a) its concentration cannot be measured or controlled \quad
(b) \ce{N2O2} \quad (c) rate $= k[\ce{A}][\ce{C}]$ \quad (d) it needs a
three-body collision}
\keyonly{\begin{apcanswer}(a) Because an intermediate is formed and
consumed inside the mechanism, so you \textbf{cannot measure or set its
concentration}. A rate law has to be usable, which means it may contain
only reactants, products or catalysts. (b) \textbf{\ce{N2O2}} --- made
in step 1, consumed in step 2, absent from the overall equation. (c)
Slow step gives $\text{rate} = k_{2}[\ce{B}][\ce{C}]$. From the fast
equilibrium, $K = [\ce{B}]/[\ce{A}]$, so $[\ce{B}] = K[\ce{A}]$.
Substituting: $\text{rate} = \mathbf{k[\ce{A}][\ce{C}]}$ with
$k = k_{2}K$. (d) Because it would require \textbf{three particles to
collide at the same instant} with enough energy and the right
orientation --- overwhelmingly improbable. The two-step
pre-equilibrium route reproduces the same third-order rate law using
only bimolecular collisions.\end{apcanswer}}
\end{yourturn}

% =====================================================================
\section*{\sffamily Ladder 10 \textbullet\ The collision model}
\zsec{12.6}

For a reaction to occur, particles must collide --- with
\textbf{enough energy} and in the \textbf{right orientation}. Almost
every collision fails one test or the other, which is why reactions are
so much slower than collision frequencies alone would suggest.

\begin{center}
\begin{tikzpicture}[font=\sffamily\footnotesize]
  \node[font=\sffamily\small\bfseries] at (2.4,3.2) {right orientation};
  \fill[apcgray!65] (0.9,2.1) circle (0.2);
  \fill[apcgray!25] (1.45,2.1) circle (0.15);
  \draw[-{Stealth[length=2mm]},thick] (2.0,2.1) -- (2.6,2.1);
  \fill[apcgray!25] (3.3,2.1) circle (0.15);
  \fill[apcgray!65] (3.85,2.1) circle (0.2);
  \node[align=center] at (2.4,1.2) {reactive ends meet\\
    \textbf{reaction happens}};

  \begin{scope}[xshift=6.0cm]
    \node[font=\sffamily\small\bfseries] at (2.4,3.2)
      {wrong orientation};
    \fill[apcgray!65] (0.9,2.1) circle (0.2);
    \fill[apcgray!25] (1.45,2.1) circle (0.15);
    \draw[-{Stealth[length=2mm]},thick] (2.0,2.1) -- (2.6,2.1);
    \fill[apcgray!65] (3.3,2.1) circle (0.2);
    \fill[apcgray!25] (3.85,2.1) circle (0.15);
    \node[align=center] at (2.4,1.2) {wrong ends meet\\
      \textbf{they simply bounce}};
  \end{scope}
\end{tikzpicture}
\end{center}

\begin{center}
\begin{tikzpicture}
\begin{axis}[
  width=11.2cm, height=6.2cm,
  xlabel={kinetic energy}, ylabel={fraction of particles},
  xmin=0, xmax=120, ymin=0, ymax=3.4,
  xtick=\empty, ytick=\empty,
  axis lines=left,
  legend pos=north east, legend style={font=\footnotesize,
    cells={anchor=west}},
  label style={font=\footnotesize},
]
\addplot[seriesA, thick, domain=0.05:120, samples=110]
  {100*sqrt(x)*exp(-x/15)/58.09};
\addlegendentry{lower temperature}
\addplot[seriesB, thick, domain=0.05:120, samples=110]
  {100*sqrt(x)*exp(-x/28)/148.2};
\addlegendentry{higher temperature}
\draw[very thick, densely dashed] (axis cs:60,0) -- (axis cs:60,2.6);
\node[font=\footnotesize,anchor=south] at (axis cs:60,2.62)
  {$E_{\mathrm{a}}$};
\node[font=\footnotesize,anchor=west] at (axis cs:63,1.15)
  {only particles out \emph{here}};
\node[font=\footnotesize,anchor=west] at (axis cs:63,0.85)
  {can react at all};
\end{axis}
\end{tikzpicture}
\end{center}

\begin{apcnote}[Why a small temperature rise does so much]
Heating does \emph{not} shift the whole distribution up --- it flattens
and stretches it to the right. The peak barely moves, but the
\textbf{tail beyond $E_{\mathrm{a}}$} grows enormously, and it is only
the tail that reacts. For a typical activation energy near
\SI{50}{\kilo\joule\per\mole}, warming from \SI{25}{\celsius} to
\SI{35}{\celsius} roughly \emph{doubles} the rate --- a
\SI{10}{\celsius} change producing a 100\% change in speed. Say
``the fraction of collisions with energy $\ge E_{\mathrm{a}}$
increases'', not ``the particles have more energy'': the second earns
nothing.
\end{apcnote}

\begin{workedexample}[I do: the Arrhenius equation]
\[ k = A\,e^{-E_{\mathrm{a}}/RT}
   \qquad\text{or, taking logs,}\qquad
   \ln k = -\frac{E_{\mathrm{a}}}{R}\cdot\frac{1}{T} + \ln A \]

\textbf{Read the second form as $y = mx + c$.} Plotting $\ln k$ against
$1/T$ gives a straight line of slope $-E_{\mathrm{a}}/R$ --- which is
how activation energies are actually measured.

\begin{center}
\begin{tikzpicture}
\begin{axis}[
  width=8.4cm, height=4.4cm,
  xlabel={$1/T$ (\si{\per\kelvin})}, ylabel={$\ln k$},
  xmin=0.0029, xmax=0.0035, ymin=-0.4, ymax=3.2,
  xtick={0.0030,0.0032,0.0034},
  scaled x ticks=false, x tick label style={/pgf/number format/fixed,
    /pgf/number format/precision=4, font=\footnotesize},
  ytick={0,1,2,3},
  grid=major, grid style={apcgray!25},
  label style={font=\footnotesize}, tick label style={font=\footnotesize},
]
\addplot[seriesA, thick, mark=*, mark size=1.5pt]
  coordinates {(0.0030,2.66) (0.0032,1.37) (0.0034,0.08)};
\node[font=\footnotesize,anchor=west] at (axis cs:0.00305,0.6)
  {slope $= -E_{\mathrm{a}}/R$};
\end{axis}
\end{tikzpicture}
\end{center}

\textbf{Worked: the rate constant doubles when the temperature rises
from \SI{300}{\kelvin} to \SI{310}{\kelvin}. Find $E_{\mathrm{a}}$.}

Use the two-point form:
\[ \ln\!\left(\frac{k_{2}}{k_{1}}\right)
   = -\frac{E_{\mathrm{a}}}{R}
     \left(\frac{1}{T_{2}} - \frac{1}{T_{1}}\right) \]
\[ \ln 2 = -\frac{E_{\mathrm{a}}}{8.314}
   \left(\frac{1}{310} - \frac{1}{300}\right)
   = -\frac{E_{\mathrm{a}}}{8.314}\left(\num{-1.075e-4}\right) \]
\[ 0.693 = E_{\mathrm{a}} \times \num{1.293e-5}
   \quad\Longrightarrow\quad
   E_{\mathrm{a}} = \SI{53600}{\joule\per\mole}
   = \mathbf{\SI{53.6}{\kilo\joule\per\mole}} \]

\textbf{Two habits.} Temperature must be in \textbf{kelvin} --- the
equation is meaningless in celsius. And $R = \SI{8.314}{}$ gives
$E_{\mathrm{a}}$ in joules, so divide by 1000 at the end; forgetting to
is the single most common slip here.
\end{workedexample}

\begin{apcnote}[AP scope: the calculation above is enrichment]
The CED's topic 5.6 exclusion: \emph{``Calculations involving the
Arrhenius equation will not be assessed on the AP Exam.''} What
\textbf{is} assessed: the collision-model reasoning, the
Maxwell--Boltzmann tail argument, and \emph{interpreting} a $\ln k$
versus $1/T$ plot --- that it is linear, that its slope is
$-E_{\mathrm{a}}/R$, and that a steeper line means a larger activation
energy. The two-point calculation is kept because Zumdahl teaches it
and it cements the relationship; on the exam, the qualitative version
carries the marks.
\end{apcnote}

\begin{yourturn}[YOUR TURN 10 --- four questions]
\begin{enumerate}[leftmargin=*,itemsep=0.5em]
  \item State the two requirements for a collision to lead to reaction.
        \workspace[1.6cm]
  \item Why does raising the temperature increase the rate so sharply?
        \workspace[1.8cm]
  \item What does the slope of a plot of $\ln k$ against $1/T$ equal?
        \workspace[1.6cm]
  \item Does raising the temperature change $E_{\mathrm{a}}$? Explain.
        \workspace[1.8cm]
\end{enumerate}
\selfcheck{(a) enough energy, correct orientation \quad (b) the
fraction exceeding $E_{\mathrm{a}}$ rises steeply \quad (c)
$-E_{\mathrm{a}}/R$ \quad (d) no --- it changes how many particles clear
it}
\keyonly{\begin{apcanswer}(a) The colliding particles must have
\textbf{kinetic energy of at least $E_{\mathrm{a}}$}, and they must meet
in the \textbf{correct orientation}. Failing either means they simply
bounce apart. (b) Because heating stretches the
Maxwell--Boltzmann distribution to the right, and the
\textbf{fraction of particles with energy $\ge E_{\mathrm{a}}$} --- the
area under the tail --- grows very steeply. A modest temperature rise
produces a large change in that fraction, which is why roughly
\SI{10}{\celsius} can double a rate. (c)
$\mathbf{-E_{\mathrm{a}}/R}$. (d) \textbf{No.} $E_{\mathrm{a}}$ is a
property of the reaction pathway itself and is unchanged by temperature.
What changes is the \emph{fraction of collisions that clear the
barrier}. Only a catalyst, which supplies a different pathway, lowers
$E_{\mathrm{a}}$.\end{apcanswer}}
\end{yourturn}

% =====================================================================
\section*{\sffamily Ladder 11 \textbullet\ Reaction energy profiles}
\zsec{12.6}

An energy profile is the reaction drawn as a landscape: reactants on one
side, products on the other, and a hill in between whose height is the
activation energy.

\begin{center}
\begin{tikzpicture}[font=\sffamily\footnotesize]
  \draw[-{Stealth[length=2.5mm]},thick] (0.5,0.3) -- (0.5,4.3)
    node[above,font=\footnotesize] {energy};
  \draw[-{Stealth[length=2.5mm]},thick] (0.5,0.3) -- (9.6,0.3)
    node[right,font=\footnotesize] {reaction progress};
  % curve: reactants 2.2, peak 3.7, products 1.2
  \draw[very thick] (1.1,2.2) -- (2.2,2.2)
    .. controls (3.3,2.2) and (3.5,3.7) .. (4.4,3.7)
    .. controls (5.3,3.7) and (5.5,1.2) .. (6.6,1.2) -- (8.6,1.2);
  % levels
  \draw[dotted,thick] (1.1,2.2) -- (8.9,2.2);
  \draw[dotted,thick] (1.1,1.2) -- (8.9,1.2);
  \draw[dotted,thick] (2.6,3.7) -- (8.9,3.7);
  % Ea forward
  \draw[{Stealth[length=2mm]}-{Stealth[length=2mm]},thick]
    (7.4,2.2) -- (7.4,3.7);
  \node[anchor=west,font=\footnotesize] at (7.5,2.95)
    {$E_{\mathrm{a}}$ forward};
  % Ea reverse
  \draw[{Stealth[length=2mm]}-{Stealth[length=2mm]},thick]
    (8.75,1.2) -- (8.75,3.7);
  \node[anchor=west,font=\footnotesize] at (8.85,2.45)
    {$E_{\mathrm{a}}$ rev.};
  % dH
  \draw[{Stealth[length=2mm]}-{Stealth[length=2mm]},thick]
    (1.55,2.2) -- (1.55,1.2);
  \node[anchor=east,font=\footnotesize] at (1.45,1.7) {$\Delta H$};
  \node[anchor=south,font=\footnotesize] at (1.6,2.25) {reactants};
  \node[anchor=north,font=\footnotesize] at (5.9,1.15) {products};
  \node[anchor=south,font=\footnotesize] at (4.4,3.75)
    {transition state};
\end{tikzpicture}
\end{center}

\[ E_{\mathrm{a}}(\text{forward}) - E_{\mathrm{a}}(\text{reverse})
   = \Delta H \]

\begin{center}
\begin{tikzpicture}[font=\sffamily\footnotesize]
  \node[font=\sffamily\small\bfseries] at (5.0,4.35)
    {a two-step mechanism has \textbf{two} humps};
  \draw[-{Stealth[length=2.5mm]},thick] (0.5,0.3) -- (0.5,4.0);
  \draw[-{Stealth[length=2.5mm]},thick] (0.5,0.3) -- (9.4,0.3)
    node[right,font=\footnotesize] {progress};
  % reactants 1.6, peak1 3.4, intermediate 2.1, peak2 2.9, products 1.0
  \draw[very thick] (1.0,1.6) -- (1.7,1.6)
    .. controls (2.5,1.6) and (2.6,3.4) .. (3.3,3.4)
    .. controls (4.0,3.4) and (4.1,2.1) .. (4.8,2.1)
    .. controls (5.5,2.1) and (5.6,2.9) .. (6.3,2.9)
    .. controls (7.0,2.9) and (7.1,1.0) .. (7.9,1.0) -- (8.8,1.0);
  \node[anchor=south,font=\footnotesize] at (3.3,3.45) {peak 1};
  \node[anchor=south,font=\footnotesize] at (6.3,2.95) {peak 2};
  \node[anchor=north,font=\footnotesize] at (4.8,2.05) {intermediate};
  \node[anchor=south,font=\footnotesize] at (1.3,1.65) {reactants};
  \node[anchor=north,font=\footnotesize] at (8.3,0.95) {products};
  % Moved to a caption line below the axis: at y=3.75 this ran straight
  % through the "peak 1" label.
  \node[font=\sffamily\footnotesize] at (4.9,-0.35)
    {the \textbf{higher} barrier is the rate-determining step};
\end{tikzpicture}
\end{center}

\begin{workedexample}[I do: reading numbers off a profile]
\textbf{A one-step reaction has $E_{\mathrm{a}}(\text{forward}) =
\SI{80}{\kilo\joule\per\mole}$ and $\Delta H =
\SI{-30}{\kilo\joule\per\mole}$. Find $E_{\mathrm{a}}$ for the reverse
reaction.}

\textbf{Put the reactants at zero and work up.} The peak is
\SI{80}{\kilo\joule\per\mole} above the reactants. The products are
\SI{30}{\kilo\joule\per\mole} \emph{below} them, because $\Delta H$ is
negative.

\textbf{The reverse barrier is measured from the products to the same
peak:}
\[ E_{\mathrm{a}}(\text{rev}) = 80 - (-30)
   = \mathbf{\SI{110}{\kilo\joule\per\mole}} \]

\textbf{Sanity check.} An exothermic reaction always has a
\emph{larger} reverse barrier than forward one --- the products sit in a
deeper valley, so climbing back out costs more. If your reverse
$E_{\mathrm{a}}$ came out smaller for an exothermic reaction, you
subtracted the wrong way.

\textbf{Now the two-step profile above.} Take reactants at 0, peak 1 at
75, the intermediate at 20, peak 2 at 50, products at $-25$
(\si{\kilo\joule\per\mole}).
\begin{itemize}[leftmargin=1.4em,itemsep=0.15em]
  \item $E_{\mathrm{a}}$ of step 1 $= 75 - 0 = \SI{75}{}$
  \item $E_{\mathrm{a}}$ of step 2 $= 50 - 20 = \SI{30}{}$
  \item overall $\Delta H = -25 - 0 = \SI{-25}{\kilo\joule\per\mole}$
\end{itemize}
\textbf{Step 1 is rate-determining} because its barrier is the larger.
Each step's activation energy is measured from \emph{its own} starting
valley, not from the reactants --- for step 2 that means from the
intermediate at 20, not from 0.
\end{workedexample}

\begin{yourturn}[YOUR TURN 11 --- four questions]
Use reactants $=0$, peak 1 $=\SI{60}{}$, intermediate $=\SI{15}{}$,
peak 2 $=\SI{45}{}$, products $=\SI{-20}{\kilo\joule\per\mole}$.
\begin{enumerate}[leftmargin=*,itemsep=0.5em]
  \item Find $E_{\mathrm{a}}$ for step 1. \workspace[1.4cm]
  \item Find $E_{\mathrm{a}}$ for step 2. \workspace[1.4cm]
  \item Which step is rate-determining? \workspace[1.4cm]
  \item Give the overall $\Delta H$, and say whether the reaction is
        exothermic or endothermic. \workspace[1.8cm]
\end{enumerate}
\selfcheck{(a) \SI{60}{\kilo\joule\per\mole} \quad (b)
\SI{30}{\kilo\joule\per\mole} \quad (c) step 1 \quad (d)
\SI{-20}{\kilo\joule\per\mole}, exothermic}
\keyonly{\begin{apcanswer}(a) $60 - 0 =
\mathbf{\SI{60}{\kilo\joule\per\mole}}$. (b) $45 - 15 =
\mathbf{\SI{30}{\kilo\joule\per\mole}}$ --- measured from the
intermediate, not from the reactants. (c) \textbf{Step 1}, because its
activation energy (\SI{60}{}) is the larger of the two, so it is the
bottleneck. (d) $-20 - 0 =
\mathbf{\SI{-20}{\kilo\joule\per\mole}}$, and a negative $\Delta H$
means the reaction is \textbf{exothermic} --- the products lie below the
reactants.\end{apcanswer}}
\end{yourturn}

% =====================================================================
\section*{\sffamily Ladder 12 \textbullet\ Catalysis}
\zsec{12.7}

A catalyst speeds a reaction by offering a \textbf{different pathway}
with a lower activation energy. It changes the route, not the
destination.

\begin{center}
\begin{tikzpicture}[font=\sffamily\footnotesize]
  \draw[-{Stealth[length=2.5mm]},thick] (0.5,0.3) -- (0.5,4.2)
    node[above,font=\footnotesize] {energy};
  \draw[-{Stealth[length=2.5mm]},thick] (0.5,0.3) -- (9.3,0.3)
    node[right,font=\footnotesize] {progress};
  % uncatalysed: reactants 1.9, peak 3.6, products 1.0
  \draw[seriesA, very thick] (1.0,1.9) -- (1.9,1.9)
    .. controls (2.8,1.9) and (3.0,3.6) .. (3.9,3.6)
    .. controls (4.8,3.6) and (5.0,1.0) .. (6.0,1.0) -- (8.4,1.0);
  % catalysed: same start and end, lower peak
  \draw[seriesB, very thick] (1.0,1.9) -- (1.9,1.9)
    % peak lowered 2.7 -> 2.55 to open a real gap for its own label,
    % which the flat top of this curve was grazing.
    .. controls (2.8,1.9) and (3.0,2.55) .. (3.9,2.55)
    .. controls (4.8,2.55) and (5.0,1.0) .. (6.0,1.0) -- (8.4,1.0);
  \node[anchor=south,font=\footnotesize] at (3.9,3.65) {uncatalysed};
  % Above its own peak, not below it: anchored north at 2.62 the
  % descending catalysed curve cut through the end of the word.
  \node[anchor=south,font=\footnotesize] at (3.9,2.68) {catalysed};
  \draw[dotted,thick] (1.0,1.9) -- (8.9,1.9);
  \draw[dotted,thick] (1.0,1.0) -- (8.9,1.0);
  \draw[{Stealth[length=2mm]}-{Stealth[length=2mm]},thick]
    (8.7,1.0) -- (8.7,1.9);
  \node[anchor=west,font=\footnotesize] at (8.8,1.45) {$\Delta H$};
  % Caption line below the axis: at y=3.9 the "uncatalysed" curve label
  % (anchored south at y=3.65) printed through it.
  \node[font=\sffamily\footnotesize] at (4.9,-0.35)
    {same start, same finish, \textbf{same $\Delta H$} --- only the hill
     is lower};
\end{tikzpicture}
\end{center}

\begin{apctrap}
\textbf{A catalyst does not change $\Delta H$, and it does not shift the
equilibrium position.} It lowers the barrier for the forward
\emph{and} the reverse reaction by exactly the same amount, so both
speed up equally and equilibrium is simply reached sooner. An answer
saying a catalyst ``increases the yield'' or ``makes the reaction more
exothermic'' scores zero.
\end{apctrap}

\begin{workedexample}[I do: catalyst or intermediate?]
\textbf{Both appear inside a mechanism and neither appears in the
overall equation. Telling them apart is a standard exam question, and
the test is the order in which they show up.}

\textbf{An intermediate is produced first, then consumed.} It appears as
a \emph{product} of an earlier step and a \emph{reactant} of a later
one. It did not exist before the reaction started.

\textbf{A catalyst is consumed first, then regenerated.} It appears as a
\emph{reactant} of an earlier step and a \emph{product} of a later one.
It was there at the start and is still there at the end --- which is why
a tiny amount can process an unlimited quantity of reactant.

\textbf{Worked case.} In the mechanism
\[ \ce{Cl + O3 -> ClO + O2}
   \qquad
   \ce{ClO + O -> Cl + O2} \]
\ce{Cl} is used in step 1 and \textbf{regenerated} in step 2 --- a
\textbf{catalyst}. \ce{ClO} is made in step 1 and \textbf{consumed} in
step 2 --- an \textbf{intermediate}. Adding the steps gives
$\ce{O3 + O -> 2O2}$, with neither appearing. This is the real
chemistry of chlorine-catalysed ozone destruction, and the reason a
single chlorine atom can destroy many thousands of ozone molecules
before it is finally removed.

\textbf{Two kinds of catalyst, distinguished by phase.}
\textbf{Homogeneous} catalysts are in the same phase as the reactants
--- as \ce{Cl} above, a gas among gases. \textbf{Heterogeneous}
catalysts are in a different phase, usually a solid surface on which gas
or liquid reactants adsorb, react and then leave; the platinum in a car's
catalytic converter is the standard example. \textbf{Enzymes} are
biological catalysts, and their extraordinary specificity comes from an
active site shaped to hold one substrate in exactly the right
orientation --- solving the orientation half of Ladder 10's problem
rather than the energy half.
\end{workedexample}

\begin{yourturn}[YOUR TURN 12 --- four questions]
\begin{enumerate}[leftmargin=*,itemsep=0.5em]
  \item How does a catalyst increase the rate? \workspace[1.6cm]
  \item Does a catalyst change $\Delta H$? Explain.
        \workspace[1.8cm]
  \item Distinguish an intermediate from a catalyst.
        \workspace[1.8cm]
  \item Does a catalyst shift the position of equilibrium? Explain.
        \workspace[1.8cm]
\end{enumerate}
\selfcheck{(a) it lowers $E_{\mathrm{a}}$ via a new pathway \quad (b)
no \quad (c) intermediate is made then used; catalyst is used then
remade \quad (d) no --- both directions speed up equally}
\keyonly{\begin{apcanswer}(a) By providing an \textbf{alternative
pathway with a lower activation energy}, so a much larger fraction of
collisions has enough energy to react. It does not give the particles
more energy. (b) \textbf{No.} $\Delta H$ depends only on the energies of
the reactants and products, and a catalyst changes neither --- it alters
only the height of the barrier between them. (c) An
\textbf{intermediate} is \emph{produced} in an earlier step and
\emph{consumed} in a later one; a \textbf{catalyst} is \emph{consumed}
in an earlier step and \emph{regenerated} in a later one. Neither
appears in the overall equation, but the catalyst exists before the
reaction begins and survives it. (d) \textbf{No.} A catalyst lowers the
activation energy of the forward and reverse reactions by the
\textbf{same amount}, so both rates increase by the same factor.
Equilibrium is reached \emph{sooner} but sits in exactly the same
place.\end{apcanswer}}
\end{yourturn}

% =====================================================================
\section*{\sffamily Mastery tracker}

Tick a row only if \textbf{all four} YOUR TURN questions were right on
the first attempt. Every row is assessed except the Arrhenius
\emph{calculation} inside Ladder 10 (see its scope note) --- the
collision model and the graph reading are fully fair game.

\renewcommand{\arraystretch}{1.3}
\noindent\begin{tabular}{@{}c p{5.9cm} c p{4.9cm}@{}}
\toprule
\textbf{First try?} & \textbf{Skill} & \textbf{Ladder} &
  \textbf{If not, re-read\ldots}\\
\midrule
$\square$ & rate and stoichiometry & 1 & divide by the coefficient \\
$\square$ & rate laws and order & 2 & orders are experimental \\
$\square$ & method of initial rates & 3 & change one thing at a time \\
$\square$ & units of $k$ & 4 & $\mathrm{M^{\,1-n}\,s^{-1}}$ \\
$\square$ & integrated rate laws & 5 & which plot is straight? \\
$\square$ & half-life & 6 & only first order is constant \\
$\square$ & elementary steps & 7 & coefficients work only here \\
$\square$ & mechanisms and the RDS & 8 & the slow step sets the rate \\
$\square$ & pre-equilibrium & 9 & eliminate the intermediate \\
$\square$ & the collision model & 10 & energy \emph{and} orientation \\
$\square$ & energy profiles & 11 & measure from its own valley \\
$\square$ & catalysis & 12 & new pathway, same $\Delta H$ \\
\bottomrule
\end{tabular}

\begin{apcnote}[Scoring yourself honestly]
12/12: you have Unit 5. This chapter is unusually self-contained ---
almost nothing in it depends on the rest of the course --- so a clean
sweep here is genuinely banked.

9--11: look at \emph{where} the misses fall. Ladders 1--6 are the
computational half and repair with practice. Ladders 7--9 are the
mechanism half, and a miss there is nearly always the same root cause:
still reading orders off coefficients. Re-read Ladder 7 before
attempting more problems.

8 or fewer: go back to Ladder 2 and hold one sentence in mind ---
\emph{orders come from experiment, not from the balanced equation}. Most
wrong answers in this unit are that error wearing a different hat.
Ladder 3 (initial rates) and Ladder 8 (mechanisms) are the two that
appear most often on the exam; if time is short, make those two
automatic first.
\end{apcnote}

\end{document}
